\section{Related work}
\label{sec:prior-art}

No existing system realizes this proposal end-to-end, but every
component has been built somewhere, and the most reliable cost
evidence for each comes from its builders. This section surveys four
bodies of work: foundational semantics for \code{unsafe} Rust, deductive verifiers
for Rust, the ``specification language in comments'' tradition, and formal
models of managed runtimes. \Cref{tab:tools} summarizes the Rust-side
tools against the axes that matter here.

\subsection{Foundational semantics}
\label{sec:prior-foundations}

RustBelt~\cite{rustbelt2018} established the template this whole proposal
rests on: interpret each type as a separation-logic predicate in
Iris~\cite{iris2018}, prove safe code well-typed once and for all, and
admit each \code{unsafe} library by \emph{manually proving its
implementation inhabits the semantic interpretation of its API}. That is
the formal content of a \safetyc{} comment. The qualification is
equally instructive: RustBelt's proofs concern
$\lambda_{\mathrm{Rust}}$, a core calculus, not surface Rust, and each
verified library (\code{Cell}, \code{Mutex}, \code{Arc}, \dots) cost
an expert person-months of work in Rocq.

Two further facts about the foundations constrain any design:

\begin{itemize}[leftmargin=1.4em]
\item \textbf{The aliasing model is part of the obligation.} A proof that
  an \code{unsafe} block is memory-safe but violates the compiler's
  aliasing discipline proves the wrong theorem: rustc optimizes under
  Stacked Borrows~\cite{stackedborrows2019} and, increasingly, Tree
  Borrows~\cite{treeborrows2025} (PLDI 2025 Distinguished Paper, now
  formalized in Rocq and implemented in Miri, and already the mode
  \dotnetrs{}'s CI runs). Any DSL that cannot state
  retag/permission obligations silently under-approximates UB.
\item \textbf{Rust has no complete normative formal semantics in 2026.}
  The Ferrocene specification was adopted as the seed of the official
  Rust specification in March 2025~\cite{fls2025}, but it is prose;
  \code{a-mir-formality}~\cite{amirformality} covers static semantics
  only; MiniRust~\cite{minirust} is the leading candidate for the
  operational story and is explicitly the UB list a proof should target.
  Every ``real theorem about the Rust substrate'' is therefore stated
  against a best-available approximation, and the study's soundness
  claims must say so.
\end{itemize}

\subsection{Deductive verifiers for Rust}
\label{sec:prior-verifiers}

\begin{table}
  \centering
  \small
  \begin{tabular}{@{}lllll@{}}
    \toprule
    Tool & Handles \code{unsafe}? & Foundation & Automation & Annotation site \\
    \midrule
    RustBelt~\cite{rustbelt2018}       & yes ($\lambda_{\mathrm{Rust}}$)  & Iris/Rocq, foundational & none & external Rocq \\
    RefinedRust~\cite{refinedrust2024} & yes (real MIR)     & Iris/Rocq, foundational & high (Lithium) & \code{\#[rr::...]} attributes \\
    VeriFast~\cite{verifastrust2022}   & yes                & SL, tool-trusted$^\dagger$ & symbolic exec. & \code{/*@...@*/} comments \\
    Verus~\cite{verus2024}             & re-encoded$^\ddagger$ & SMT; Iris justification~\cite{verusbelt2026} & very high (Z3) & \code{verus!\{\}} macro \\
    Creusot~\cite{creusot2022}         & emerging           & Why3/Coma & high (portfolio) & attributes \\
    Gillian-Rust~\cite{gillianrust2025}& yes (hybrid)       & SL, prototype & high & attributes \\
    Aeneas~\cite{aeneas2022}           & \textbf{no}        & Lean 4 (pure translation) & Lean tactics & external Lean \\
    Prusti                             & no                 & Viper & high & attributes \\
    Flux~\cite{flux2023}               & no                 & liquid types/SMT & very high & attributes \\
    Kani~\cite{kani2026}               & bounded only       & CBMC & total & harnesses \\
    \bottomrule
  \end{tabular}
  \caption{Rust verification tools relevant to the proposal.
  $^\dagger$Foundational certification of VeriFast results is in progress
  via hinted mirroring~\cite{foundationalverifast2026}.
  $^\ddagger$Verus does not verify preexisting \code{unsafe} blocks; code
  is rewritten against trusted linear-permission APIs
  (\code{PointsTo}, \code{PPtr}, \code{PCell}).}
  \label{tab:tools}
\end{table}

Three entries warrant elaboration because they delimit the design
space.

\paragraph{RefinedRust is the proposal, minus the ECMA-335 layer.}
RefinedRust~\cite{refinedrust2024} attaches refinement-type annotations to
real Rust source, translates the actual MIR into a Rocq model (Radium),
and discharges obligations with untrusted-but-foundationally-checked
automation inside Iris; it verified a \code{Vec} variant including
intricate raw-pointer code. It demonstrates that annotation-driven,
machine-checked, foundational proofs about real \code{unsafe} Rust
exist today. It also supplies the clearest cost evidence: it was built
by a team of Iris experts, language coverage remains partial, and
difficult cases still require manual Iris lemmas. Wholesale adoption
by a single developer would itself be a research project; adopting its
\emph{architecture} --- annotations, a mechanized model, and
foundationally checked automation --- at narrower scope is the
realistic variant.

\paragraph{VeriFast is the precedent for proofs that live in comments.}
VeriFast for Rust~\cite{verifastrust2022,verifasttb2026} puts contracts,
predicate definitions, and proof steps (lemma calls, \code{open}/%
\code{close}) in \code{/*@ ... @*/} comments beside the code, verifies
modules containing \code{unsafe} blocks by symbolic execution, derives
contracts for safe functions from RustBelt's type interpretation, and ---
critically for this proposal --- recent work adds Tree Borrows
compliance obligations.
It is one of the tools in the AWS \code{verify-rust-std}
campaign~\cite{verifyruststd2025} and verified properties of std's
\code{LinkedList}. Its trust story (an unverified checker core, with
foundational certification arriving via hinted
mirroring~\cite{foundationalverifast2026}) is exactly the trust position a
custom \safetyc{} checker would occupy --- a point \cref{sec:design} leans
on.

\paragraph{Verus marks the automation ceiling.}
Verus~\cite{verus2024,verusghost2023} demonstrates what SMT combined
with linear ghost permissions can achieve: systems code with pointer
manipulation verified at industrial scale (Amazon, Microsoft), with
proofs written in Rust syntax --- and, since
VerusBelt~\cite{verusbelt2026}, an Iris/Rocq semantic soundness proof
backing its permission axioms. Its limitation for \dotnetrs{} is
structural: Verus is designed to own the code under verification. One
does not annotate an existing \gcarena{}-based heap; one rewrites it
against \code{PPtr}/\code{PCell} tokens. Verus is also the ecosystem
with the most evidence for LLM-generated proofs (AutoVerus, SAFE),
which is relevant to the intended workflow.

Kani~\cite{kani2026} and Miri sit below all of these: bounded model
checking and dynamic interpretation respectively. They prove no theorems,
but they falsify candidate invariants cheaply, and the
\code{verify-rust-std} campaign~\cite{verifyruststd2025} (16{,}748
auto-generated harnesses; contract attributes now upstreamed into nightly
std as an experiment) is the best available evidence on what annotating
\code{unsafe} Rust at scale costs in practice.

\subsection{Structured safety comments: an emerging convention}
\label{sec:prior-tags}

The Rust community is independently converging on machine-readable
\safetyc{} comments. RFC 3842 ``Safety Tags''~\cite{rfc3842} (opened July
2025, still open) proposes \code{\#[safety::requires(tag, ...)]} on
\code{unsafe} functions and \code{\#[safety::checked(tag, ...)]} at call
sites; the \code{tag-std} project~\cite{tagstd2025} developed the
underlying taxonomy and a linter that checks \emph{structural consistency}
--- every obligation named by the callee is acknowledged by the caller ---
and applied it to libstd, covering 96.1\% of public unsafe APIs. The
authors note explicitly that tags ``could support formal verification,
e.g., by being extended into contracts.''

These systems check bookkeeping, not truth: a tag asserts that someone
claims alignment holds, not that it does. That is precisely the gap a
proof layer fills, and it fixes the surface syntax question: the DSL
should be an \emph{extension of} the emerging tag convention (tag $=$
theorem name; proof term or proof-script reference attached per tag), not
a fork of it. The older traditions --- ACSL/Frama-C~\cite{acsl} on C and
JML/OpenJML~\cite{openjml} on Java --- are the existence proof for
comment-carried, machine-checked specifications at industrial scale, and
they teach two lessons: the WP-style split (SMT for the bulk, proof
assistant for the residue) is what makes the burden bearable, and
comment-hosted languages drift from the host language unless the
compiler sees them --- which favors attribute/macro embedding for anything
load-bearing.

\subsection{Formal models of managed runtimes}
\label{sec:prior-clr}

The specification-side picture is starkly asymmetric to the Rust side:
\textbf{no mechanized formalization of ECMA-335 exists.} What exists is
pre-mechanization-era and pen-and-paper:

\begin{itemize}[leftmargin=1.4em]
\item \textbf{Fruja's ASM model}~\cite{fruja2006,borger2005csharp} --- the
  most complete formal CLR model ever built: modular Abstract State
  Machines covering the CIL instruction set, the exception mechanism
  (including filter/finally/fault interactions \dotnetrs{} implements as
  a state machine), and the bytecode verifier, with type-safety proofs.
  Dormant since $\sim$2007, never mechanized. As a \emph{source text} for
  a mechanization it is invaluable: the semantic edge cases are already
  worked out.
\item \textbf{Gordon and Syme}~\cite{gordonsyme2001} --- type soundness
  for ``Baby IL,'' the theory behind the real verifier, defining the
  evaluation-stack typing and managed-pointer discipline that several of
  \dotnetrs{}'s invariant families instantiate.
\item \textbf{Yu, Kennedy, and Syme}~\cite{kennedysyme2004,kennedysyme2001}
  --- soundness for the CLR's generics implementation strategy (shared
  instantiations, runtime type representations); the invariant catalog
  for any theorem about \dotnetrs{}'s generic-instantiation caches.
\item \textbf{Benton}~\cite{benton2005} --- a compositional program logic
  for a CIL-like stack machine; the early precedent for Hoare-style
  reasoning at the evaluation-stack level.
\end{itemize}

The JVM side shows what mechanization costs and what it buys.
Jinja~\cite{jinja2006} and JinjaThreads~\cite{jinjathreads2012} grew, over
roughly fifteen years of Isabelle/HOL work, into a verified multithreaded
Java-like machine with a verified compiler and executable semantics ---
proof that a full managed-bytecode machine is mechanizable, and a warning
that it is among the largest artifacts in its proof assistant's library.
JSCert~\cite{jscert2014} contributes the methodology this study adopts for
the ECMA-335 model: clause-by-clause correspondence between the prose
standard and the mechanized definitions, plus an \emph{executable}
reference interpreter differentially tested against real test suites ---
a philosophy \dotnetrs{} already practices with its fixture harness
against the real .NET runtime.

For the GC specifically, the verified-collector literature supplies the
invariant vocabulary. Hawblitzel and Petrank~\cite{hawblitzelpetrank2009}
verified practical copying and mark-sweep collectors with Boogie/Z3 ---
two people, SMT-automated, performance-competitive --- and their
mutator--collector interface contracts (reachability abstraction,
forwarding invariants) transfer directly to the obligations \dotnetrs{}
discharges by hand around \gcarena{}. CakeML's verified generational
collector~\cite{cakemlgc2019} and the rely-guarantee verification of a
concurrent collector~\cite{zakowski2018} mark the high end;
Verve~\cite{verve2010} shows a whole runtime under a typed managed
language verified with SMT-heavy tooling. The relevant reading for
\dotnetrs{} is mutator-side: the project does not implement a collector,
it \emph{uses} one, so its obligations are correct usage of \gcarena{}'s
rooting/tracing contract --- which is precisely the shape of the
Hawblitzel--Petrank interface contracts.

\subsection{Proof-assistant backends}
\label{sec:prior-backends}

The two candidate backends the author named are both viable, with sharply
different prior art:

\paragraph{Rocq.} Renamed from Coq at version 9.0 (March
2025)~\cite{rocq}. Every foundational unsafe-Rust result lives here:
Iris~\cite{iris2018}, RustBelt, the Tree Borrows mechanization,
RefinedRust, VerusBelt. If theorems must genuinely bottom out in the Rust
aliasing model, Rocq is the only backend where that model already exists
mechanized. Translation tools exist (\code{rocq-of-rust}~\cite{rocqofrust}
via THIR; \code{hax}~\cite{hax} for a panic-free safe subset).

\paragraph{Lean 4.} The stronger \emph{engineering} ecosystem circa
2026~\cite{lean4}: best-in-class metaprogramming for building a DSL
elaborator, rapidly improving automation (\code{omega}, the \code{grind}
tactic since v4.22, \code{lean-smt}), Mathlib as background library, and
by far the largest LLM-proof corpus --- directly relevant to this
project's agent-driven workflow. Aeneas~\cite{aeneas2022} targets Lean as
its primary backend (safe code only). An active Iris port
exists (\code{iris-lean}~\cite{irislean}: MoSeL, \code{UPred},
invariants, later credits), but it is a port of the \emph{logic} --- there
is no mechanized Rust semantics in Lean today.

The resulting design fork, developed in \cref{sec:design}: Rocq for
the Rust-substrate half, Lean for the ECMA-335-model half, or one
backend with the other half stated against a trusted axiomatization.
